\documentclass[a5paper,10pt]{article}

% https://github.com/InDevRus/filippov-solutions

% Нумерация страниц
\pagenumbering{gobble}

% Настройка полей
\usepackage{geometry}
\geometry{left=1cm}
\geometry{right=1cm}
\geometry{top=1cm}
\geometry{bottom=1.5cm}

% Вставка таблицы
\usepackage{array}
\usepackage{tabularx}

% Инструменты выравнивания формул
\usepackage{amsmath}
\usepackage{mathtools}

% Диакритика в математике
\usepackage{amsfonts}

% Вычисления размеров на ходу
\usepackage{calc}

% Удобные записи производных
\usepackage[italic=true]{derivative}

% Настройки сносок
\usepackage{enotez}

% Длинные знаки векторов
\usepackage{anyfontsize}
\usepackage[b]{esvect}

% Вставка картинок
\usepackage{graphicx}

% Вставка красной строки
\usepackage{indentfirst}
\setlength{\parindent}{1em}

% Условия в командах
\usepackage{xifthen}

% Луа-скрипты
\usepackage{luacode}

% Настройка команд отдельно в математическом и текстовом режимах
\usepackage{mathcommand}

% Для повторения LaTeX кода
\usepackage{multido}

% Конфигурация отступов формул
\delimitershortfall-1sp
\usepackage{mleftright}
\mleftright

% Растягивание объектов
\usepackage{relsize}

% Обтекание текстом
\usepackage{wrapfig}
\usepackage{subcaption}
\usepackage{floatflt}

% Поддержка ввода кириллицы
\usepackage[english, russian]{babel}

% Настройка корневой позиции для компилятора
\usepackage{import}
\providecommand*{\ProjectRootPath}{..}

\import{\ProjectRootPath/preambles/macros/}{theorems}

\import{\ProjectRootPath/preambles/fonts}{font_setup}

% Знак градуса
\usepackage{siunitx}

\import{\ProjectRootPath/preambles/macros/}{operators}
\import{\ProjectRootPath/preambles/macros/}{redefinitions}

\import{\ProjectRootPath/preambles/fonts}{letters_setup}
\import{\ProjectRootPath/preambles/fonts/adjustments}{subsuperscripts}

\import{\ProjectRootPath/preambles/custom}{frames}
\import{\ProjectRootPath/preambles/custom}{formatting}
\import{\ProjectRootPath/preambles/custom}{list_setup}

% TODO: перевести команды на современный стиль объявления

\begin{document}
    $\br{u - x} \parder{u} {x} + \br{u - y} \parder{u} {y} - z \parder{u} {y} = x + y$.
    
    $\dvfrac{\di x} {u - x} = \dvfrac{\di y} {u - y} = \dvfrac{\di z} {-z} = \dvfrac{\di u} {x + y}$.

    \begin{enumerate}
        \item $\dvfrac{\di x - \di y} {y - x} = \dvfrac{\di z} {-z}$.
        $\ln\vbr{\y - x} - \ln\vbr{z} = A$.
        $\dvfrac{y - x}{z} = A$.

        \item $\dvfrac{\di x + \di y + 2 \di u} {2u + x + y} = \dvfrac{\di z} {-z}$.
        $\ln\vbr{2u + x + y} + \ln\vbr{z} = B$.
        $\br{2u + x + y}z = B$.

        \item $\dvfrac{\di x + \di y} {2u - x - y} = \dvfrac{\di u} {x + y}$.
        Переобозначим $x + y = v$.
        $\dvfrac{\di v} {2u - v} = \dvfrac{\di u} {v}$.
        Это однородное уравнение 1-ой степени, особая точка которого является <<седлом>>. 
        Подбором находим два частных решения (сепаратрисы) $v - u = 0$ и $v + 2u = 0$.
        Замена переменных
        $$\tsys{
            \phi = \dfrac{v - u}{3} \vspace{0.5em} \\
            \xi = \dfrac{v + 2u}{3},
        }
        \text{ т.е. }
        \tsys{
            v = 2\phi + \xi \\
            u = \xi - \phi
        }$$
        приводит к $\xi \di \phi + 2\phi \di \xi = 0$,
        откуда $\phi \pow{\xi}{2} = C$; $\br{x + y - u} \br{x + y + 2u}^2 = C$.
    \end{enumerate} 

    Запись ответа можно упростить, заменив $A$ на
    $D = A \cdot B = \br{\y - x}\br{2u + x + y}$. Итак,
    $F\br{\br{\y - x}\br{2u + x + y}; \br{2u + x + y}z; \br{x + y - u} \br{x + y + 2u}^2} = 0$.
\end{document}
